% GENERATED FILE. Edit canonical lessons, then run npm run generate:books.
% canonical-source-hash: fnv1a64:9c060a5dda90fd91
% canonical-lessons: RU-C373-riskovat, RU-C373-opasatsya, RU-C373-izbegat, RU-C373-ugrozhat, RU-C373-predotvrashchat

\chapter{Playing it safe}
\label{ch:ru-playing-it-safe}
\hlchaptermodality{373}

\begin{chapteropening}
\textbf{By the end of this chapter:} \emph{I can talk about taking a risk, being wary of something, avoiding it, saying that something threatens us, and preventing it.}

Complete the last lesson of chapter 373: I can talk about taking a risk, being wary of something, avoiding it, saying that something threatens us, and preventing it.
\end{chapteropening}

\section[\texorpdfstring{\ru{рисковать} (riskovát)}{рисковать (riskovát)}]{\ru{рисковать} (riskovát) — to take a risk, to risk}

\label{lesson:RU-C373-riskovat}



\begin{quote}
Before the new one: say the Russian for a chance, then the Russian for a risk.
\end{quote}

\subsection*{You'll want to know: \ru{рисковать}}
\textbf{\ru{рисковать}} (\emph{riskovát}) — \textquotedblleft{}to take a risk, to risk\textquotedblright{}.

\begin{quote}
\ru{Я} \ru{не} \ru{хочу} \ru{рисковать} \ru{своими} \ru{деньгами}.
\end{quote}

\begin{quote}
\emph{Ya ne khochú riskovát svoími dengámi.}
\end{quote}

\begin{quote}
I do not want to risk my money.
\end{quote}

Imperfective; the present is \emph{\ru{рискую}, \ru{рискуешь}}; perfective \emph{\ru{рискнуть}}. \emph{\ru{Рисковать}} + instrumental — to risk something.

\begin{cousinweb}
\textbf{\ru{рисковать}} \textquotedblleft{}to risk\textquotedblright{} (EtymDB) is built on \emph{\ru{риск}}, which is usually traced, through French \emph{risque}, to Italian \emph{risco}.
\end{cousinweb}

\begin{grammarlens}[title={What you've built}]
One more word for telling, explaining and describing at B1.
\end{grammarlens}

\subsection*{Your turn}
\begin{itemize}
\raggedright
  \item \emph{Say it:} \emph{riskovát}
  \item \emph{Say it:} \emph{riskovát}, once more
  \item \emph{Say it:} a sentence of your own with \emph{riskovát}
  \item \emph{Recall:} say the Russian for a chance, then the Russian for a risk, then say \emph{riskovát} again
\end{itemize}

\begin{tcolorbox}[breakable,colback=teal!4,colframe=teal!35!black,title={Before you move on}]
What does \ru{рисковать} mean? (\textquotedblleft{}To take a risk, to risk\textquotedblright{}.) Say it once more.
\end{tcolorbox}
\section[\texorpdfstring{\ru{опасаться} (opasátsya)}{опасаться (opasátsya)}]{\ru{опасаться} (opasátsya) — to fear, to be wary of}

\label{lesson:RU-C373-opasatsya}



\begin{quote}
Before the new one: say the Russian for a risk, then the Russian for to take a risk.
\end{quote}

\subsection*{You'll want to know: \ru{опасаться}}
\textbf{\ru{опасаться}} (\emph{opasátsya}) — \textquotedblleft{}to fear, to be wary of\textquotedblright{}.

\begin{quote}
\ru{Врачи} \ru{опасаются} \ru{новой} \ru{эпидемии}.
\end{quote}

\begin{quote}
\emph{Vrachí opasáyutsya nóvoy epidémii.}
\end{quote}

\begin{quote}
Doctors fear a new epidemic.
\end{quote}

Imperfective only. It takes the genitive: \emph{\ru{опасаться} \ru{дождя}}. More formal than \emph{\ru{бояться}}.

\begin{cousinweb}
\textbf{\ru{опасаться}} is usually traced to \emph{\ru{о}-} and \emph{\ru{пасти}} \textquotedblleft{}to guard, to watch over\textquotedblright{}, inherited from Proto-Slavic \emph{*pasti} (EtymDB): to be wary is to keep watch. \emph{\ru{Опасный}} \textquotedblleft{}dangerous\textquotedblright{} is built on it.
\end{cousinweb}

\begin{grammarlens}[title={What you've built}]
One more word for telling, explaining and describing at B1.
\end{grammarlens}

\subsection*{Your turn}
\begin{itemize}
\raggedright
  \item \emph{Say it:} \emph{opasátsya}
  \item \emph{Say it:} \emph{opasátsya}, once more
  \item \emph{Say it:} a sentence of your own with \emph{opasátsya}
  \item \emph{Recall:} say the Russian for a risk, then the Russian for to take a risk, then say \emph{opasátsya} again
\end{itemize}

\begin{tcolorbox}[breakable,colback=teal!4,colframe=teal!35!black,title={Before you move on}]
What does \ru{опасаться} mean? (\textquotedblleft{}To fear, to be wary of\textquotedblright{}.) Say it once more.
\end{tcolorbox}
\section[\texorpdfstring{\ru{избегать} (izbegát)}{избегать (izbegát)}]{\ru{избегать} (izbegát) — to avoid}

\label{lesson:RU-C373-izbegat}



\begin{quote}
Before the new one: say the Russian for to take a risk, then the Russian for to fear.
\end{quote}

\subsection*{You'll want to know: \ru{избегать}}
\textbf{\ru{избегать}} (\emph{izbegát}) — \textquotedblleft{}to avoid\textquotedblright{}.

\begin{quote}
\ru{Он} \ru{избегает} \ru{трудных} \ru{разговоров}.
\end{quote}

\begin{quote}
\emph{On izbegáyet trúdnykh razgovórov.}
\end{quote}

\begin{quote}
He avoids difficult conversations.
\end{quote}

Imperfective; perfective \emph{\ru{избежать}}, conjugated like \emph{\ru{бежать}}. It takes the genitive: \emph{\ru{избегать} \ru{ошибок}}.

\begin{cousinweb}
\textbf{\ru{избегать}} is built from \emph{\ru{из}-} \textquotedblleft{}out of\textquotedblright{} and \emph{\ru{бегать}} \textquotedblleft{}to run\textquotedblright{}, from Proto-Slavic \emph{*běžati}: to run out of a thing’s way.
\end{cousinweb}

\begin{grammarlens}[title={What you've built}]
One more word for telling, explaining and describing at B1.
\end{grammarlens}

\subsection*{Your turn}
\begin{itemize}
\raggedright
  \item \emph{Say it:} \emph{izbegát}
  \item \emph{Say it:} \emph{izbegát}, once more
  \item \emph{Say it:} a sentence of your own with \emph{izbegát}
  \item \emph{Recall:} say the Russian for to take a risk, then the Russian for to fear, then say \emph{izbegát} again
\end{itemize}

\begin{tcolorbox}[breakable,colback=teal!4,colframe=teal!35!black,title={Before you move on}]
What does \ru{избегать} mean? (\textquotedblleft{}To avoid\textquotedblright{}.) Say it once more.
\end{tcolorbox}
\section[\texorpdfstring{\ru{угрожать} (ugrozhát)}{угрожать (ugrozhát)}]{\ru{угрожать} (ugrozhát) — to threaten}

\label{lesson:RU-C373-ugrozhat}



\begin{quote}
Before the new one: say the Russian for to fear, then the Russian for to avoid.
\end{quote}

\subsection*{You'll want to know: \ru{угрожать}}
\textbf{\ru{угрожать}} (\emph{ugrozhát}) — \textquotedblleft{}to threaten\textquotedblright{}.

\begin{quote}
\ru{Сильный} \ru{ветер} \ru{угрожает} \ru{старым} \ru{деревьям}.
\end{quote}

\begin{quote}
\emph{Sílnyy véter ugrozháyet stárym derévyam.}
\end{quote}

\begin{quote}
The strong wind threatens the old trees.
\end{quote}

Imperfective. \emph{\ru{Угрожать} \ru{кому} \ru{чем}}: the person threatened goes in the dative. The noun is \emph{\ru{угроза}}, a threat.

\begin{cousinweb}
\textbf{\ru{угрожать}} is built from \emph{\ru{у}-} and \emph{\ru{гроза}} \textquotedblleft{}a thunderstorm\textquotedblright{}, which in older Russian also meant terror and threat; \emph{\ru{з}} softens to \emph{\ru{ж}} in the verb.
\end{cousinweb}

\begin{grammarlens}[title={What you've built}]
One more word for telling, explaining and describing at B1.
\end{grammarlens}

\subsection*{Your turn}
\begin{itemize}
\raggedright
  \item \emph{Say it:} \emph{ugrozhát}
  \item \emph{Say it:} \emph{ugrozhát}, once more
  \item \emph{Say it:} a sentence of your own with \emph{ugrozhát}
  \item \emph{Recall:} say the Russian for to fear, then the Russian for to avoid, then say \emph{ugrozhát} again
\end{itemize}

\begin{tcolorbox}[breakable,colback=teal!4,colframe=teal!35!black,title={Before you move on}]
What does \ru{угрожать} mean? (\textquotedblleft{}To threaten\textquotedblright{}.) Say it once more.
\end{tcolorbox}
\section[\texorpdfstring{\ru{предотвращать} (predotvrashchát)}{предотвращать (predotvrashchát)}]{\ru{предотвращать} (predotvrashchát) — to prevent}

\label{lesson:RU-C373-predotvrashchat}



\begin{quote}
Before the new one: say the Russian for to avoid, then the Russian for to threaten.
\end{quote}

\subsection*{You'll want to know: \ru{предотвращать}}
\textbf{\ru{предотвращать}} (\emph{predotvrashchát}) — \textquotedblleft{}to prevent\textquotedblright{}.

\begin{quote}
\ru{Эти} \ru{правила} \ru{предотвращают} \ru{аварии}.
\end{quote}

\begin{quote}
\emph{Éti právila predotvrashcháyut avárii.}
\end{quote}

\begin{quote}
These rules prevent accidents.
\end{quote}

Imperfective; the perfective is \emph{\ru{предотвратить}}, with \emph{\ru{предотвращу}, \ru{предотвратишь}}. Formal.

\begin{cousinweb}
\textbf{\ru{предотвращать}} is built from \emph{\ru{пред}-} \textquotedblleft{}before\textquotedblright{}, \emph{\ru{от}-} \textquotedblleft{}away\textquotedblright{} and \emph{\ru{вращать}} \textquotedblleft{}to turn\textquotedblright{}, the Church Slavonic twin of native \emph{\ru{ворочать}}: to turn a danger away in advance. The \emph{-\ru{ра}-} and the \emph{\ru{щ}} are both Church Slavonic shapes.
\end{cousinweb}

\begin{grammarlens}[title={What you've built}]
One more word for telling, explaining and describing at B1.
\end{grammarlens}

\subsection*{Your turn}
\begin{itemize}
\raggedright
  \item \emph{Say it:} \emph{predotvrashchát}
  \item \emph{Say it:} \emph{predotvrashchát}, once more
  \item \emph{Say it:} a sentence of your own with \emph{predotvrashchát}
  \item \emph{Recall:} say the Russian for to avoid, then the Russian for to threaten, then say \emph{predotvrashchát} again
\end{itemize}

\begin{tcolorbox}[breakable,colback=teal!4,colframe=teal!35!black,title={Before you move on}]
What does \ru{предотвращать} mean? (\textquotedblleft{}To prevent\textquotedblright{}.) Say it once more.
\end{tcolorbox}
